\chapter{Návrh emulátoru}

% Výběr platformy, jestli to bude osobní počítač s operačním systémem nebo pouhý mikrokontrolér, přímo
% ovlivňuje strukturu a návrh emulátoru. Jelikož naší cílovou platformou je právě mikrokontrolér
% RP2350, musí být emulátor navržen dostatečně lehce, ale stále odpovídat ostatním požadavkům.

% - emulator je nizkourovnovy
% - emulator je sekvencni
% - emulator je cyklove presny
% - emulator realizuje catch up
% - CPU je ridicim prvkem

Emulátor je navržen jako nízkoúrovňový, což znamená, že emuluje na úrovni jednotlivých
instrukcí procesoru, registrů a cyklů. Tento přístup zajišťuje lepší kompatibilitu s původním
softwarem, který často využíval nedokumentované vlastnosti hardwaru NES.

Vzhledem k tomu, že RP2350 disponuje omezenými prostředky pro paralelní zpracování s vysokou režií
(např. multitasking v reálném čase), je emulátor navržen jako sekvenční a využívá catch-up metodu
(viz \ref{tp:catch-up}) pro synchronizaci komponentů. Paralelní běh by vyžadoval náročnou
synchronizaci a zamykání prostředků, takže tento přístup zjednodušuje návrh a je optimální
na RP2350.

CPU je řídicí prvek: PPU a APU ,,spí'', dokud CPU nevykoná určitý počet cyklů nebo
nepokusí se s nimi komunikovat. V ten moment tyto komponenty dohánějí stav aktuální stav CPU. Díky
tomuto modelu je emulátor zároveň cyklově přesný. Každá instrukce CPU trvá přesně stejný počet
emulovaných taktů jako na reálném hardwaru, což je kritické pro hry, které využívají přesné časování
pro grafické efekty (např. změna palety uprostřed vykreslování řádku, nebo zápis do registru
scrollu).

Pro reprezentaci komponentů v kódu se používají struktury, které zapouzdřují celý vnitřní stav
komponentů. Struktury neobsahuji pouze registry, ale i vnitřní čítače, příznaky a případně buffery.
To umožňuje v budoucnu například snadné ukládání stavu systému do souboru a jeho opětovné načtení.
Ke každé struktuře náleží sada funkcí, zpravidla zodpovědné za inicializaci, resetování,
aktualizaci stavu a sběrnicové funkce (\texttt{read()}/\texttt{write()}).

\subsection{Smyčka emulátoru}

Smyčka emulátoru provádí emulaci po snímcích PPU (viz výpis \ref{lst:emu_loop}). Nejdřív se určí
počet cyklů, které vykoná CPU, a až pak se aktualizuje stav PPU a APU (dohánějí CPU). Pokud při
dohánění na vyžádání PPU dokončí vykreslování snímku, CPU musí ihned ukončit vykonání, aby nedošlo k
desynchronizaci.

\begin{listing}[htbp]
\begin{minted}{c}
/* Emulace se provádí po snímcích */
while (!ppu->frame_done) {
    CycleCounter cycles_to_run = /* ... */;
    /* CPU provede určitý počet cyklů */
    while (cpu->cycles < cycles_to_run && !ppu->frame_done) {
        cpu_step(cpu);
    }
    /* Dohánění */
    ppu_run_until(ppu, cpu->cycles * 3);
    apu_run_until(apu, cpu->cycles);
}

ppu->frame_done = false;
\end{minted}
\caption{Smyčka emulátoru. CPU provádí určitý počet instrukcí, a pak PPU a APU dohánějí jeho stav.}
\label{lst:emu_loop}
\end{listing}

Potřebný počet cyklů pro CPU se určuje na základě nejbližší významné událostí -- stav VBlank,
přerušení NMI nebo konec vykreslovaného řádku. Pro NTSC verzi systému NES platí že za jeden řádek
CPU vykoná cca~113,67 cyklů.

\section{ROM obrazy a mappery}

Abstrakce herní kazety a její vnitřní logiky je realizována skrze strukturu \texttt{Disk}. Tato
struktura zapouzdřuje veškerá metadata o načteném ROM obrazu, ukazatele na alokované paměťové bloky
(PRG a CHR) a aktuální vnitřní stav aktivního mapperu.

Inicializace struktury probíhá pomocí funkcí \texttt{disk\_load()} (pro načítání ze souborového
systému SD karty) a \texttt{disk\_load\_mem()} (pro přímé zavedení z paměti RAM). Obě funkce
provádějí parsování a validaci hlavičky formátu iNES, inicializaci mapperu a alokaci PRG a CHR
pamětí na základě zjištěných údajů.

Vzhledem k tomu, že každý mapper vyžaduje pro svou činnost odlišné řídicí proměnné -- jako jsou
čítače přerušení, registry bankování nebo latch stavy -- bylo nutné vyřešit jejich ukládání. Pro
úsporu paměti je stav mapperu v rámci struktury \texttt{Disk} definován jako sjednocení
(\texttt{union}). Tento přístup zajišťuje, že v paměti je vyhrazen pouze takový prostor, který
odpovídá nejnáročnějšímu podporovanému mapperu. Jednotlivé stavy se díky vlastnostem typu
\texttt{union} v paměti překrývají, což znemožňuje jejich současnou existenci (která v emulátoru
nikdy nenastává) a výrazně snižuje celkovou paměťovou stopu objektu \texttt{Disk}.

V aktuální fázi vývoje emulátor implementuje podporu tři klíčových mapperů, které pokrývají
podstatnou část knihovny her pro systém NES: NROM, UNROM a MMC3.
